% ===========================================================================
%  ch_preamble.tex --- Chapter M (Mathematical introduction).
%  Adapted from 7source_claude_1/ch2_preamble.tex.  When this chapter is
%  dropped into the thesis, merge this file into the thesis preamble;
%  duplicates may simply be deleted.
% ===========================================================================

\usepackage[T1]{fontenc}
\usepackage[utf8]{inputenc}
\usepackage{lmodern}
\usepackage{microtype}
\usepackage{amsmath,amssymb,amsthm,mathtools}
\usepackage{graphicx}
\usepackage{float}
\usepackage{placeins}
\usepackage{booktabs}
\usepackage[hypcap=false]{caption}
\usepackage{subcaption}
\usepackage{tikz}
\usepackage{pgfplots}
\pgfplotsset{compat=1.18}
\usepackage{enumitem}
\usepackage{xcolor}
\usepackage{array}
\usepackage{hyperref}
\usepackage[nameinlink,noabbrev]{cleveref}

% --- Equation cross-references as bare parenthesised numbers -----------------
\crefformat{equation}{(#2#1#3)}
\Crefformat{equation}{(#2#1#3)}
\crefrangeformat{equation}{(#3#1#4)--(#5#2#6)}
\crefmultiformat{equation}{(#2#1#3)}{ and (#2#1#3)}{, (#2#1#3)}{ and (#2#1#3)}

% --- Shared style for the cobweb panels of the Galton--Watson section ---------
\pgfplotsset{
    ch2plotstyle/.style={
        width=4.9cm, height=4.9cm,
        axis lines=left,
        xlabel={$S$},
        xmin=0, xmax=1, ymin=0, ymax=1,
        xtick={0,0.5,1}, ytick={0,0.5,1},
        grid=major, grid style={dashed, gray!30},
        samples=100, domain=0:1,
        thick,
        tick label style={font=\tiny},
        label style={font=\small}
    }
}

% --- Notation ---------------------------------------------------------------
\newcommand{\pgf}{probability generating function}
\newcommand{\rv}{random variable}
\newcommand{\pr}[1]{\mathbb{P}\!\left(#1\right)}
\newcommand{\ex}[1]{\mathbb{E}\!\left(#1\right)}
\newcommand{\var}[1]{\operatorname{Var}\!\left(#1\right)}
\newcommand{\dd}{\mathrm{d}}
\newcommand{\ee}{\mathrm{e}}
\newcommand{\ddt}[1]{\frac{\dd #1}{\dd t}}
\newcommand{\dda}[2]{\frac{\dd #1}{\dd #2}}
\newcommand{\pddt}[1]{\frac{\partial #1}{\partial t}}
\newcommand{\pdd}[2]{\frac{\partial #1}{\partial #2}}
\newcommand{\lrb}[1]{\left(#1\right)}
\newcommand{\lrbs}[1]{\left[#1\right]}
\newcommand{\Ac}{A_{\mathrm{c}}}
\newcommand{\ind}{\mathbf{1}}

% --- Cross-chapter references ------------------------------------------------
% Standalone-build form, numbered for the expected thesis order (mathematical
% introduction = Chapter 2, the constant A(p) = Chapter 3).  Thesis integration
% will redefine \ChA to a numbered reference (e.g.\ Chapter~\ref{ch:Ap}).
\newcommand{\ChA}{Chapter~3 (\emph{The constant $A(p)$})}

% --- Forward references ------------------------------------------------------
% Every forward reference to later thesis chapters is routed through \fwd,
% whose first argument is a stable key and whose second is the prose actually
% typeset.  Once the thesis chapter order is fixed, redefining \fwd alone is
% enough to make every forward reference carry a chapter number.
\newcommand{\fwd}[2]{#2}

% --- Theorem environments ----------------------------------------------------
\theoremstyle{plain}
\newtheorem{theorem}{Theorem}[chapter]
\newtheorem{proposition}[theorem]{Proposition}
\newtheorem{lemma}[theorem]{Lemma}
\newtheorem{corollary}[theorem]{Corollary}
\theoremstyle{definition}
\newtheorem{definition}[theorem]{Definition}
\theoremstyle{remark}
\newtheorem{remark}[theorem]{Remark}

\allowdisplaybreaks
